% Canonical derivation source for the Green-function decay crossover.
% This file is a LaTeX fragment and is included by the Supplementary Information.

\section{Bessel--Morse asymptotics of the Green response}
\label{sec:green-function-asymptotics}

We derive the time-domain response of the ZGV branch with every Fourier,
modal, and stationary-phase constant retained.  The object propagated below
is the positive-frequency complex response.  It is not doubled.  The
physical real response is reconstructed only after the branch integral has
been evaluated.  This distinction fixes a factor of two that is otherwise
easy to hide in a loosely defined analytic signal.

Throughout this section \(a>0\), \(k_0>0\), and the time dependence of a
positive-frequency mode is \(\exp(-\mathrm{i}\omega t)\).  The hypotheses
on branch simplicity, eigengap, and smoothness are those of
Theorem~\ref{thm:anisotropic-normal-form}.  All asymptotic statements concern
the selected positive-frequency branch in a fixed annular spectral window;
contributions from other branches are separate terms in the complete Green
function.

\subsection{Fourier convention and normalization-invariant branch projection}

For an in-plane field, use the transform pair
\begin{equation}
  \widehat{\mathbf u}(\mathbf k,z,t)
  =\int_{\mathbb R^2}\mathbf u(\mathbf x,z,t)
    \mathrm e^{-\mathrm i\mathbf k\cdot\mathbf x}\,\mathrm d^2\mathbf x,
  \qquad
  \mathbf u(\mathbf x,z,t)
  =\mathscr F_2\int_{\mathbb R^2}
    \widehat{\mathbf u}(\mathbf k,z,t)
    \mathrm e^{\mathrm i\mathbf k\cdot\mathbf x}\,\mathrm d^2\mathbf k,
  \qquad
  \mathscr F_2=(2\pi)^{-2}.
  \label{eq:green-fourier-convention}
\end{equation}
At fixed \(\mathbf k\), let the thickness-discretized or weak-form dynamics
under an impulsive generalized load be
\begin{equation}
  M\ddot{\mathbf U}+K\mathbf U
  =\mathbf f\,\delta(t),
  \qquad
  K\mathbf u=\omega^2M\mathbf u,
  \qquad
  N:=\mathbf u^\dagger M\mathbf u>0.
  \label{eq:impulsive-hermitian-dynamics}
\end{equation}
The vector \(\mathbf f=\mathbf f(\mathbf k)\) includes the through-thickness
source functional, and a linear detector is written
\(D[\mathbf U]=\mathbf d^\dagger\mathbf U\).  Projecting
Eq.~\eqref{eq:impulsive-hermitian-dynamics} onto a simple mode gives
\begin{equation}
  \ddot q+\omega^2q
  =\frac{\mathbf u^\dagger\mathbf f}{N}\,\delta(t),
  \qquad
  q(t)=H(t)\frac{\mathbf u^\dagger\mathbf f}{\omega N}\sin(\omega t).
  \label{eq:projected-impulse-oscillator}
\end{equation}
Because
\begin{equation}
  \frac{\sin(\omega t)}{\omega}
  =2\operatorname{Re}\left[
      \frac{\mathrm i}{2\omega}\mathrm e^{-\mathrm i\omega t}
    \right],
  \label{eq:positive-frequency-oscillator-split}
\end{equation}
the undoubled coefficient of the positive-frequency exponential in the
detected response is
\begin{equation}
  \mathcal P(\mathbf k)
  :=\frac{\mathrm i}{2\omega}
    \frac{(\mathbf d^\dagger\mathbf u)
          (\mathbf u^\dagger\mathbf f)}
         {\mathbf u^\dagger M\mathbf u}.
  \label{eq:source-mode-detector-amplitude}
\end{equation}
This source--mode--detector amplitude is invariant under an arbitrary
nonzero complex normalization \(\mathbf u\mapsto c(\mathbf k)\mathbf u\):
the numerator and denominator both acquire \(|c|^2\), including when the
scale and phase depend on \(\mathbf k\).  Thus no mass-normalization
convention enters the observable amplitude.

In the mass-normalized, colocated top-face normal channel,
\(\mathbf d=\mathbf f=\mathbf e_z(h)\) and
\(\mathbf u^\dagger M\mathbf u=1\).  The general projection then reduces to
the phase-invariant overlap used by the numerical implementation,
\begin{equation}
  \mathcal P_n(\mathbf k)
  =\frac{\mathrm i}{2\omega_n}
    \left|u_{n,z}(h;\mathbf k)\right|^2.
  \label{eq:normal-impulse-overlap}
\end{equation}

Let \(G^{+}(t)\) denote the inverse Fourier integral made from the coefficient
in Eq.~\eqref{eq:source-mode-detector-amplitude}.  For a real load and
a real detector, the negative-frequency term is its complex conjugate after
the reciprocal wave-vector partner is included.  Consequently
\begin{equation}
  G_{\mathrm{phys}}(t)
  =G^{+}(t)+G^{+}(t)^*
  =2\operatorname{Re}G^{+}(t).
  \label{eq:physical-real-reconstruction}
\end{equation}
Equation~\eqref{eq:physical-real-reconstruction}, not an extra factor inside
\(G^{+}\), defines the physical real response.  Every theorem below is first
stated for \(G^{+}\), and the physical result follows from this equation.

For reference, a normalized circular Gaussian source of radius \(r_s\),
\begin{equation}
  s(\mathbf x)=\frac{1}{\pi r_s^2}
    \exp\left(-\frac{|\mathbf x|^2}{r_s^2}\right),
  \qquad
  \widehat s(k)=\exp\left[-\left(\frac{r_sk}{2}\right)^2\right],
  \label{eq:gaussian-source-transform}
\end{equation}
multiplies \(\mathcal P\) by \(\widehat s(k)\).  The numerical local-branch
window is the smooth super-Gaussian multiplied by the compact taper used by
the direct numerical integral,
\begin{equation}
  \chi_{\sigma_q}(q)=\exp[-(q/\sigma_q)^8]b_{\sigma_q}(q),
  \qquad q=k-k_0,
  \label{eq:green-super-gaussian-window}
\end{equation}
although an ordinary Gaussian window
\(\exp[-q^2/(2\sigma_q^2)]\) gives the same leading asymptotics.  Only
smoothness, localization inside the annulus, and a nonzero value at the
stationary set are used in the proof.  Here \(b_{\sigma_q}=1\) on
\(|q|\leq1.25\sigma_q\), \(b_{\sigma_q}=0\) on
\(|q|\geq1.5\sigma_q\), and the intervening strip uses the standard flat
step constructed from \(\exp(-1/x)\).  Hence \(b_{\sigma_q}\) is
\(C^\infty\), all endpoint derivatives vanish, every local coefficient below
is unchanged, and the numerical integral satisfies the same no-endpoint
hypothesis as the theorem.

Define the complete projected amplitude, excluding the Fourier measure and
the polar Jacobian, by
\begin{equation}
  \mathcal A(q,\theta;\varepsilon)
  :=\chi_{\sigma_q}(q)
    \exp\left[-\left(\frac{r_s(k_0+q)}{2}\right)^2\right]
    \frac{\mathrm i}{2\omega}
    \frac{(\mathbf d^\dagger\mathbf u)
          (\mathbf u^\dagger\mathbf f_z)}
         {\mathbf u^\dagger M\mathbf u}.
  \label{eq:complete-green-amplitude}
\end{equation}
Here \(\mathbf f_z\) is the unit through-thickness source functional; any
additional calibrated source strength can be included multiplicatively.
For the colocated axisymmetric normal channel in the isotropic plate,
\begin{equation}
  \mathcal A(0,\theta;0)=A_0,
  \qquad A_0\ne0.
  \label{eq:canonical-nonnodal-amplitude}
\end{equation}
In particular, \(A_0\) does not contain \(k_0\), \(\mathscr F_2\), or the
angular factor \(2\pi\).

\subsection{Local polar integral and radial stationary phase}

Let \(0<\eta<k_0\) be small enough that the selected branch remains simple
on \(|q|\leq\eta\), and let the window vanish smoothly before the annular
boundary if a compact cutoff is needed.  The exact local
positive-frequency complex response is
\begin{equation}
  G^{+}(t;\varepsilon)
  =\mathscr F_2
    \int_0^{2\pi}\!\int_{-\eta}^{\eta}
    \mathcal A(q,\theta;\varepsilon)
    \exp[-\mathrm it\omega(q,\theta;\varepsilon)]
    (k_0+q)\,\mathrm dq\,\mathrm d\theta.
  \label{eq:branch-projected-response}
\end{equation}
This displays the complete Fourier/polar measure
\((2\pi)^{-2}(k_0+q)\,\mathrm dq\,\mathrm d\theta\).

For sufficiently small \(|\varepsilon|\),
Theorem~\ref{thm:anisotropic-normal-form} supplies the unique radial
stationary graph \(q=q_*(\theta,\varepsilon)\).  Uniform one-dimensional
stationary phase in \(q\), with the angular variable treated as a compact
parameter, gives
\begin{align}
  G^{+}(t;\varepsilon)
  ={}&\mathscr F_2\,
    \mathrm e^{-\mathrm i\pi/4}
    \left(\frac{2\pi}{t}\right)^{1/2}
    \int_0^{2\pi}
    \frac{(k_0+q_*)\mathcal A(q_*,\theta;\varepsilon)}
         {\sqrt{\omega_{qq}(q_*,\theta;\varepsilon)}}
    \mathrm e^{-\mathrm it\omega_*(\theta;\varepsilon)}
    \,\mathrm d\theta
    +O(t^{-3/2}),
  \label{eq:radial-stationary-reduction}\\
  \omega_*(\theta;\varepsilon)
  ={}&\omega_0+\varepsilon V(\theta)
    +\varepsilon^2W_2^{\mathrm{eff}}(\theta)
    +O_{C^2}(\varepsilon^3),                                      \notag\\
  W_2^{\mathrm{eff}}(\theta)
  :={}&W_2(\theta)-\frac{B(\theta)^2}{2a}.
  \label{eq:effective-second-order-green-phase}
\end{align}
The last line is exactly the radial-relaxation correction derived in
Eq.~\eqref{eq:reduced-angular-frequency-second-order}.  Keeping \(W_2\)
but dropping \(-B^2/(2a)\) would give an incorrect second-order phase and,
therefore, an incorrect overlap window.

The radial onset time follows by requiring the stationary width
\((at)^{-1/2}\) to be small compared with the window width \(\sigma_q\):
\begin{equation}
  t_{\mathrm{rad}}\asymp\frac{1}{a\sigma_q^2}.
  \label{eq:radial-stationary-onset-time}
\end{equation}
The Fresnel phase in Eq.~\eqref{eq:radial-stationary-reduction} follows from
the exact continuation
\begin{equation}
  \int_{-\infty}^{\infty}
    \exp\left(-\frac{\mathrm i}{2}atq^2\right)\,\mathrm dq
  =\mathrm e^{-\mathrm i\pi/4}
    \sqrt{\frac{2\pi}{at}},
  \qquad a>0,\quad t>0,
  \label{eq:radial-fresnel-factor}
\end{equation}
where the oscillatory integral is defined by Gaussian continuation:
\begin{equation}
  \lim_{\eta\downarrow0}
  \int_{-\infty}^{\infty}
    \exp\left[-\left(\eta+\frac{\mathrm i}{2}at\right)q^2\right]
    \,\mathrm dq
  =\mathrm e^{-\mathrm i\pi/4}\sqrt{\frac{2\pi}{at}}.
  \label{eq:regularized-radial-fresnel-factor}
\end{equation}

\subsection{The joint transition limit and its exact Bessel kernel}

For the [001] cubic family,
\(V(\theta)=V_0+V_4\cos4\theta\), with \(V_4\) the sensitivity coefficient
and \(\varepsilon V_4\) the physical frequency shift.  Introduce
\begin{equation}
  \tau:=\varepsilon V_4t.
  \label{eq:green-transition-variable}
\end{equation}
The sign of \(\tau\) is retained; the resulting \(J_0\) kernel is even.

\begin{theorem}[Uniform Bessel crossover of the ZGV ring]
\label{thm:uniform-bessel-crossover}
Assume that the exact phase is \(C^6\), the amplitude is \(C^4\), the
radial curvature satisfies \(a>0\), \(V_4\ne0\), and
Eq.~\eqref{eq:canonical-nonnodal-amplitude} holds.  In the joint limit
\begin{equation}
  t\longrightarrow\infty,
  \qquad
  \varepsilon\longrightarrow0,
  \qquad
  \tau=\varepsilon V_4t\in K,
  \label{eq:distinguished-joint-limit}
\end{equation}
where \(K\subset\mathbb R\) is any fixed compact set,
\begin{equation}
  G^{+}(t;\varepsilon)
  =\mathcal C\,
    \mathrm e^{-\mathrm i(\omega_0+\varepsilon V_0)t}
    t^{-1/2}J_0(\varepsilon V_4t)
    +o(t^{-1/2}),
  \label{eq:uniform-response}
\end{equation}
uniformly for \(\tau\in K\), with the fully determined constant
\begin{equation}
  \mathcal C
  =\mathscr F_2(2\pi)k_0A_0
    \mathrm e^{-\mathrm i\pi/4}
    \sqrt{\frac{2\pi}{a}}.
  \label{eq:uniform-response-prefactor}
\end{equation}
No amplitude, phase, carrier-frequency correction, or time shift is fitted
in Eqs.~\eqref{eq:uniform-response}--\eqref{eq:uniform-response-prefactor}.
The corresponding physical real response is
\(G_{\mathrm{phys}}=2\operatorname{Re}G^{+}\).
\end{theorem}

\paragraph{Proof.}
Under Eq.~\eqref{eq:distinguished-joint-limit}, one has
\(\varepsilon=O(t^{-1})\).  Equations~\eqref{eq:radial-shift} and
\eqref{eq:effective-second-order-green-phase} therefore imply, uniformly in
\(\theta\),
\begin{align}
  q_*&=O(t^{-1}),
  &t\varepsilon^2W_2^{\mathrm{eff}}&=O(t^{-1}),
  &tO(\varepsilon^3)&=O(t^{-2}),                                    \notag\\
  \frac{(k_0+q_*)\mathcal A(q_*,\theta;\varepsilon)}
       {\sqrt{\omega_{qq}(q_*,\theta;\varepsilon)}}
  &=\frac{k_0A_0}{\sqrt a}+O(t^{-1}).
  \label{eq:joint-limit-radial-estimates}
\end{align}
Substitution into Eq.~\eqref{eq:radial-stationary-reduction} leaves the
angular identity
\begin{align}
  \int_0^{2\pi}\mathrm e^{-\mathrm i\tau\cos4\theta}\,\mathrm d\theta
  &=4\int_0^{\pi/2}\mathrm e^{-\mathrm i\tau\cos4\theta}\,\mathrm d\theta
    =\int_0^{2\pi}\mathrm e^{-\mathrm i\tau\cos\phi}\,\mathrm d\phi
    =2\pi J_0(\tau).
  \label{eq:angular-bessel-identity}
\end{align}
Here the change \(\phi=4\theta\) covers four identical periods.  Combining
Eq.~\eqref{eq:angular-bessel-identity} with the radial Fresnel factor gives
Eq.~\eqref{eq:uniform-response-prefactor} exactly.  The error is uniform in
absolute size on compact \(\tau\)-sets.  At a zero of \(J_0\), the displayed
leading term vanishes, so the theorem does not assert a uniform relative
error after division by \(J_0\).  \(\square\)

\subsection{General angular weights and the Fourier--Bessel selection rule}

Axisymmetry is not required for a uniform angular representation.  Let the
leading amplitude on the ring have the absolutely convergent Fourier series
\begin{equation}
  A(\theta)=\sum_{\ell\in\mathbb Z}A_\ell\mathrm e^{\mathrm i\ell\theta},
  \qquad
  A_\ell=\frac{1}{2\pi}\int_0^{2\pi}
    A(\theta)\mathrm e^{-\mathrm i\ell\theta}\,\mathrm d\theta.
  \label{eq:angular-amplitude-fourier-series}
\end{equation}
The Jacobi--Anger expansion, in the sign convention used here, is
\begin{equation}
  \mathrm e^{-\mathrm i\tau\cos4\theta}
  =\sum_{n\in\mathbb Z}(-\mathrm i)^nJ_n(\tau)
    \mathrm e^{\mathrm i4n\theta}.
  \label{eq:fourfold-jacobi-anger}
\end{equation}
Define the normalized angular kernel
\begin{equation}
  K_A(\tau):=\frac{1}{2\pi}\int_0^{2\pi}A(\theta)
    \mathrm e^{-\mathrm i\tau\cos4\theta}\,\mathrm d\theta.
  \label{eq:normalized-weighted-bessel-kernel}
\end{equation}
Termwise integration yields the exact selection rule
\begin{equation}
  K_A(\tau)
  =\sum_{n\in\mathbb Z}
    A_{-4n}(-\mathrm i)^nJ_n(\tau)
  =\sum_{\substack{\ell\in\mathbb Z\\4\mid\ell}}
    A_\ell(-\mathrm i)^{-\ell/4}J_{-\ell/4}(\tau).
  \label{eq:weighted-bessel-series}
\end{equation}
Thus only amplitude harmonics divisible by four survive the full-period
integral at this order.  The canonical law is the \(A_0J_0\) term.  If
\(A_0=0\) but another selected coefficient is nonzero, the leading uniform
kernel is the fixed linear combination of compatible higher-order Bessel
functions prescribed by Eq.~\eqref{eq:weighted-bessel-series}.  It reduces
to one higher-order Bessel function only when a single compatible harmonic
pair is present; it is not legitimate to restore a \(J_0\) term by fitting a
complex constant.
The numerical weighted-kernel routine evaluates precisely \(K_A\).  The
factor \(2\pi\) from the unnormalized angular integral already appears once,
and only once, in \(\mathcal C\) in
Eq.~\eqref{eq:uniform-response-prefactor}.

For a pure even harmonic of general order \(m\), the same change of variable
gives
\begin{equation}
  \int_0^{2\pi}\mathrm e^{-\mathrm i\tau\cos[m(\theta-\theta_0)]}
    \,\mathrm d\theta=2\pi J_0(\tau),
  \qquad m\in2\mathbb N.
  \label{eq:general-harmonic-bessel-identity}
\end{equation}
The independence of the constant from \(m\) will reappear as an exact
cancellation between Morse-point multiplicity and angular curvature.

\subsection{Ring, overlap, and fixed-anisotropy regimes}

Put
\begin{equation}
  \alpha:=\varepsilon V_4,
  \qquad
  \omega_c:=\omega_0+\varepsilon V_0,
  \qquad
  t_c:=|\alpha|^{-1}.
  \label{eq:crossover-scales}
\end{equation}
Three regimes must be kept distinct.

\paragraph{Pre-crossover ring regime.}
If
\begin{equation}
  t_{\mathrm{rad}}\ll t\ll t_c,
  \label{eq:pre-crossover-window}
\end{equation}
then \(|\tau|\ll1\) and
\begin{equation}
  J_0(\tau)=1-\frac{\tau^2}{4}+O(\tau^4).
  \label{eq:small-bessel-expansion}
\end{equation}
The continuous stationary ring has already localized radially but not
angularly, so the nonnodal response has the Morse--Bott law
\(|G^{+}|=O(t^{-1/2})\).

\paragraph{Bessel--Morse overlap.}
For \(|\tau|\gg1\), the large-argument expansion is
\begin{equation}
  J_0(\tau)
  =\sqrt{\frac{2}{\pi|\tau|}}
    \left[
      \cos\left(|\tau|-\frac{\pi}{4}\right)
      +O(|\tau|^{-1})
    \right].
  \label{eq:large-bessel-expansion}
\end{equation}
The first-order angular phase remains accurate in this late-Bessel regime
only while its accumulated higher-order error is small.  Let \(M_3\) be a
uniform bound such that the reduced-phase remainder after the displayed
second-order term is at most \(M_3|\varepsilon|^3\) in \(C^4\).  The correct
condition is
\begin{equation}
  t\gg t_{\mathrm{rad}},
  \qquad
  |\varepsilon V_4|t\gg1,
  \qquad
  \varepsilon^2t
    \lVert W_2^{\mathrm{eff}}\rVert_{L^\infty(0,2\pi)}\ll1,
  \qquad
  |\varepsilon|^3tM_3\ll1,
  \qquad
  W_2^{\mathrm{eff}}=W_2-\frac{B^2}{2a}.
  \label{eq:bessel-morse-overlap-window}
\end{equation}
When \(V_4\), \(W_2^{\mathrm{eff}}\), and the radial scales are order one,
this overlap has the familiar form
\(|\varepsilon|^{-1}\ll t\ll|\varepsilon|^{-2}\).
The compact-\(\tau\) theorem does not, by itself, justify taking
\(|\tau|\to\infty\).  The overlap requires a second, growing-parameter
stationary-phase argument.

\begin{theorem}[Growing-parameter Bessel--Morse overlap]
\label{thm:growing-bessel-morse-overlap}
Assume the hypotheses of Theorem~\ref{thm:uniform-bessel-crossover}, with
the phase and amplitude bounds uniform in a fixed small parameter interval.
Let \(\varepsilon\to0\) and \(t\to\infty\) so that, with
\(\lambda:=|\varepsilon V_4|t\),
\begin{equation}
  \frac{t}{t_{\mathrm{rad}}}\to\infty,
  \qquad
  \lambda\to\infty,
  \qquad
  \varepsilon^2t
    \lVert W_2^{\mathrm{eff}}\rVert_{L^\infty}\to0,
  \qquad
  |\varepsilon|^3tM_3\to0.
  \label{eq:growing-overlap-limit}
\end{equation}
Then
\begin{equation}
  G^{+}(t;\varepsilon)
  =\mathcal C\sqrt{\frac{2}{\pi|\alpha|}}
    \frac{\mathrm e^{-\mathrm i\omega_ct}}{t}
    \cos\left(|\alpha|t-\frac{\pi}{4}\right)
    +\mathcal E_{\mathrm{ov}}(t,\varepsilon),
  \label{eq:bessel-overlap-response}
\end{equation}
where, uniformly in this two-parameter regime,
\begin{equation}
  \mathcal E_{\mathrm{ov}}
  =\frac{1}{t\sqrt{|\alpha|}}
    O\!\left(
      \lambda^{-1}+|\varepsilon|
      +\varepsilon^2t
        \lVert W_2^{\mathrm{eff}}\rVert_{L^\infty}
      +|\varepsilon|^3tM_3
    \right)
    +O(t^{-3/2}).
  \label{eq:growing-overlap-error}
\end{equation}
Consequently Eq.~\eqref{eq:bessel-overlap-response} is a relative
asymptotic equivalence on subsequences bounded away from the cancellation
zeros, for example when
\(\left|\cos(\lambda-\pi/4)\right|\ge c>0\).
\end{theorem}

\paragraph{Proof.}
After removing the carrier in
Eq.~\eqref{eq:radial-stationary-reduction}, define
\begin{align}
  \Psi_\varepsilon(\theta)
  &:=\frac{\omega_*(\theta;\varepsilon)-\omega_c}{|\alpha|}
    =s\cos4\theta+O_{C^4}(|\varepsilon|),
    \qquad s:=\operatorname{sign}(\alpha),                         \notag\\
  \mathcal B_\varepsilon(\theta)
  &:=\frac{(k_0+q_*)\mathcal A(q_*,\theta;\varepsilon)}
          {\sqrt{\omega_{qq}(q_*,\theta;\varepsilon)}}
    =\frac{k_0A_0}{\sqrt a}+O_{C^2}(|\varepsilon|).
  \label{eq:growing-overlap-normalized-data}
\end{align}
The \(C^4\) phase estimate follows from the stated \(C^6\) regularity after
shrinking the annulus.  The normalized perturbation is order
\(\varepsilon^2/|\alpha|=|\varepsilon|/|V_4|\).  Hence the eight simple
critical points of \(s\cos4\theta\) persist, exhaust the angular circle,
and satisfy
\begin{equation}
  \theta_j(\varepsilon)=\frac{j\pi}{4}+O(|\varepsilon|),
  \qquad
  |\Psi_\varepsilon''(\theta_j)|=16+O(|\varepsilon|).
  \label{eq:growing-overlap-critical-data}
\end{equation}
One-dimensional stationary phase with large parameter \(\lambda\), applied
directly to this perturbed phase, gives
\begin{equation}
  \int_0^{2\pi}\mathcal B_\varepsilon
    \mathrm e^{-\mathrm i\lambda\Psi_\varepsilon}\,\mathrm d\theta
  =\left(\frac{2\pi}{\lambda}\right)^{1/2}
    \sum_{j=0}^{7}
    \frac{\mathcal B_\varepsilon(\theta_j)
      \mathrm e^{-\mathrm i\lambda\Psi_\varepsilon(\theta_j)
      -\mathrm i\pi\operatorname{sign}
        (\Psi_\varepsilon''(\theta_j))/4}}
         {\sqrt{|\Psi_\varepsilon''(\theta_j)|}}
    +O(\lambda^{-3/2}).
  \label{eq:growing-angular-stationary-phase}
\end{equation}
Expanding the eight coefficients back to their unperturbed values costs
\(O(|\varepsilon|)\).  Their accumulated phase correction is
\(O(\varepsilon^2t\lVert W_2^{\mathrm{eff}}\rVert_{L^\infty}
  +|\varepsilon|^3tM_3)\), as required in
Eq.~\eqref{eq:growing-overlap-limit}.  Four angular curvatures
are positive and four are negative.  Multiplying their one-dimensional
signature factors by the radial factor \(\mathrm e^{-\mathrm i\pi/4}\)
produces, respectively, the minimum factor \(-\mathrm i\) and the saddle
factor \(1\).  Summing the two fourfold families gives the leading term in
Eq.~\eqref{eq:bessel-overlap-response}; the angular remainder and the
radial \(O(t^{-3/2})\) remainder give
Eq.~\eqref{eq:growing-overlap-error}.  This proof is independent of the
compact-\(\tau\) limit. \(\square\)

Away from the cancellation zeros, the envelope is
\(O(|\alpha|^{-1/2}t^{-1})\).  Pointwise logarithmic slopes through the
zeros are undefined; an upper, RMS, or carrier-demodulated envelope must be
used with a declared fitting window.

\paragraph{Bessel constants as a Morse sum.}
The agreement is stronger than equality of exponents.  Consider the general
first-order phase
\begin{equation}
  \omega=\omega_c+\frac a2q^2
    +\alpha\cos[m(\theta-\theta_0)],
  \qquad m\in2\mathbb N,\quad \alpha\ne0.
  \label{eq:pure-harmonic-overlap-phase}
\end{equation}
There are \(m\) points at the lower frequency
\(\omega_-=\omega_c-|\alpha|\) and \(m\) at the upper frequency
\(\omega_+=\omega_c+|\alpha|\).  Since \(a>0\), the lower points are
minima and the upper points are saddles.  To leading order their Cartesian
Hessians satisfy
\begin{equation}
  \sqrt{|\det H_j|}
  =\frac{m\sqrt{a|\alpha|}}{k_0},
  \qquad
  \operatorname{sig}(H_-)=2,
  \qquad
  \operatorname{sig}(H_+)=0.
  \label{eq:overlap-hessian-signatures}
\end{equation}
The \(m\)-fold multiplicity of each family cancels the factor \(m\) in
Eq.~\eqref{eq:overlap-hessian-signatures}.  For constant amplitude \(A_0\),
ordinary two-dimensional stationary phase therefore gives
\begin{align}
  G_{\mathrm{Morse}}^{+}
  \sim{}&\frac{\mathscr F_2(2\pi)k_0A_0}
               {t\sqrt{a|\alpha|}}
    \left[
      \mathrm e^{-\mathrm i\omega_-t-\mathrm i\pi/2}
      +\mathrm e^{-\mathrm i\omega_+t}
    \right].
  \label{eq:first-order-morse-overlap-sum}
\end{align}
On the other hand, substituting Eq.~\eqref{eq:uniform-response-prefactor}
into Eq.~\eqref{eq:bessel-overlap-response} and expanding the cosine gives
\begin{align}
  G_{\mathrm{Bessel}}^{+}
  \sim{}&\frac{\mathscr F_2(2\pi)k_0A_0}
               {t\sqrt{a|\alpha|}}
  \left[
    \mathrm e^{-\mathrm i\pi/4}
    \mathrm e^{-\mathrm i\omega_ct}
    \mathrm e^{\mathrm i(|\alpha|t-\pi/4)}
    +
    \mathrm e^{-\mathrm i\pi/4}
    \mathrm e^{-\mathrm i\omega_ct}
    \mathrm e^{-\mathrm i(|\alpha|t-\pi/4)}
  \right]                                                        \notag\\
  ={}&\frac{\mathscr F_2(2\pi)k_0A_0}
               {t\sqrt{a|\alpha|}}
    \left[
      \mathrm e^{-\mathrm i\omega_-t-\mathrm i\pi/2}
      +\mathrm e^{-\mathrm i\omega_+t}
    \right].
  \label{eq:bessel-morse-overlap}
\end{align}
Thus the Fourier constant, both frequencies, the minimum and saddle
signature phases, and both multiplicities match exactly.  For the cubic
case \(m=4\), this is the equality between four minima plus four saddles and
the large-argument asymptotics of the single \(J_0\) kernel.

\subsection{Exact Cartesian Morse theorem at fixed anisotropy}

The overlap calculation still uses the first-order phase.  The ultimate
fixed-anisotropy result must instead use the exact branch frequency, exact
critical points, and exact Cartesian Hessians.

\begin{theorem}[Fixed-anisotropy Morse decay]
\label{thm:fixed-anisotropy-decay}
Fix a sufficiently small \(\varepsilon\ne0\).  Suppose the exact selected
branch \(\omega(\mathbf k;\varepsilon)\) and amplitude
\(\mathcal A(\mathbf k;\varepsilon)\) are smooth on the support of the
local window, every stationary point \(\mathbf k_j\) there is isolated and
nondegenerate, and no stationary point lies on the support boundary.  Let
\begin{equation}
  \omega_j:=\omega(\mathbf k_j;\varepsilon),
  \qquad
  H_j:=\nabla_{\mathbf k}^2\omega(\mathbf k_j;\varepsilon),
  \qquad
  \sigma_j:=\operatorname{sig}(H_j)=n_+(H_j)-n_-(H_j),
  \qquad
  \mathcal A_j:=\mathcal A(\mathbf k_j;\varepsilon).
  \label{eq:exact-morse-data}
\end{equation}
Then, as \(t\to\infty\) at this fixed \(\varepsilon\),
\begin{equation}
  G^{+}(t;\varepsilon)
  =\frac{2\pi\mathscr F_2}{t}
    \sum_j
    \frac{\mathcal A_j
      \mathrm e^{-\mathrm it\omega_j-\mathrm i\pi\sigma_j/4}}
         {\sqrt{|\det H_j|}}
    +O_\varepsilon(t^{-2}).
  \label{eq:morse-stationary-sum}
\end{equation}
Hence a generic nonnodal isolated critical point contributes \(t^{-1}\).
The constant in the remainder may diverge as \(\varepsilon\to0\), so
Eq.~\eqref{eq:morse-stationary-sum} is deliberately not a uniform
coalescing-critical-point formula.  The physical real response is again
twice the real part of the right-hand side.
\end{theorem}

\paragraph{Proof and polar-Jacobian cancellation.}
Apply the two-dimensional stationary-phase theorem in Cartesian
wave-vector coordinates to Eq.~\eqref{eq:branch-projected-response}.
For the phase \(-\omega\), the signature factor is
\(\exp[-\mathrm i\pi\operatorname{sig}(H_j)/4]\), and the Cartesian Fourier
measure is simply \(\mathscr F_2\,\mathrm d^2\mathbf k\).  This directly
gives Eq.~\eqref{eq:morse-stationary-sum}.

The same result in polar coordinates makes an important cancellation
explicit.  At a critical point, let
\begin{equation}
  H_j^{(q,\theta)}
  :=\begin{pmatrix}
      \omega_{qq}&\omega_{q\theta}\\
      \omega_{q\theta}&\omega_{\theta\theta}
    \end{pmatrix}_{(q_j,\theta_j)}.
  \label{eq:polar-coordinate-phase-hessian}
\end{equation}
Equation~\eqref{eq:critical-cartesian-polar-hessian} implies
\begin{equation}
  \det H_j^{(q,\theta)}
  =k_j^2\det H_j,
  \qquad
  \frac{k_j\mathcal A_j}
       {\sqrt{|\det H_j^{(q,\theta)}|}}
  =\frac{\mathcal A_j}{\sqrt{|\det H_j|}},
  \qquad k_j>0.
  \label{eq:polar-jacobian-cancellation}
\end{equation}
Thus the polar factor \(k_j\) cancels against the Hessian transformation.
The amplitude \(\mathcal A_j\) in the Cartesian Morse sum contains the
modal projection, Gaussian source, and branch window, but it does not
contain a residual polar Jacobian.  A minimum has signature \(+2\), a saddle
has signature \(0\), and a maximum has signature \(-2\), giving the factors
\(-\mathrm i\), \(1\), and \(+\mathrm i\), respectively.  This completes
the proof. \(\square\)

\subsection{Frequency separation and signed modulation}

For the exact fourfold branch, symmetry groups the eight points into a
fourfold minimum family and a fourfold saddle family.  Define their exact
frequencies by \(\omega_{\min}(\varepsilon)\) and
\(\omega_{\mathrm{sad}}(\varepsilon)\).  Their exact critical-frequency
separation and its weak-anisotropy expansion are
\begin{equation}
  \Delta\omega_{\mathrm{ms}}^{\mathrm{exact}}
  :=|\omega_{\mathrm{sad}}-\omega_{\min}|,
  \qquad
  \Delta\omega_{\mathrm{ms}}^{\mathrm{exact}}
  =2|\varepsilon V_4|+O(\varepsilon^2).
  \label{eq:critical-frequency-separation}
\end{equation}
Factoring a mean carrier from two equal-weight leading terms produces a
signed modulation whose angular rate is half the separation:
\begin{equation}
  \Omega_{\mathrm{mod}}
  =\frac{\Delta\omega_{\mathrm{ms}}^{\mathrm{exact}}}{2}
  =|\varepsilon V_4|+O(\varepsilon^2).
  \label{eq:signed-modulation-rate}
\end{equation}
Here \(\Omega_{\mathrm{mod}}\ge0\) is the magnitude of the angular rate of a
signed oscillatory modulation, not a signed frequency.  In the first-order
overlap its signed-cosine period and magnitude/RMS repetition period are
\begin{equation}
  T_{\mathrm{signed}}=\frac{2\pi}{\Omega_{\mathrm{mod}}},
  \qquad
  T_{|\cdot|}=\frac{\pi}{\Omega_{\mathrm{mod}}}
  =\frac{2\pi}{\Delta\omega_{\mathrm{ms}}^{\mathrm{exact}}}.
  \label{eq:modulation-periods}
\end{equation}
In the first-order overlap the signature phases make the signed factor
precisely \(\cos(|\varepsilon V_4|t-\pi/4)\), as shown in
Eq.~\eqref{eq:bessel-overlap-response}.  Its magnitude repeats twice during
one signed-cosine period.  Therefore the critical-frequency separation and
the signed modulation rate differ by a factor of two.  They should be
reported separately rather than combined under an ambiguous lossless
continuum ``beat frequency.''  At fixed \(\varepsilon\), unequal exact
amplitudes or higher-order phase shifts generally turn the pure cosine into
a two-phasor modulation without changing the exact separation.

\subsection{Nodal channels, noncommuting limits, and theorem boundaries}
\label{rem:nodal-channels}

The nonzero-amplitude hypotheses are essential rather than cosmetic.
There are three distinct nodal mechanisms.
\begin{enumerate}
  \item If the isotropic angular mean \(A_0\) vanishes, the canonical
        \(J_0\) term is absent.  Equation~\eqref{eq:weighted-bessel-series}
        supplies the correct leading Bessel order when another compatible
        Fourier coefficient is nonzero.
  \item If \(\mathcal A_j=0\) at an exact Morse point, that point has no
        \(t^{-1}\) contribution in Eq.~\eqref{eq:morse-stationary-sum}.
        The first nonzero jet of the amplitude at the point determines the
        next stationary-phase term; its exponent and phase must be derived
        for that source/detector channel.
  \item Even when every local amplitude is nonzero, symmetry-related terms
        can cancel at discrete times or in a selected detector combination.
        Such zeros do not change the generic local stationary-phase power.
\end{enumerate}
These cases also explain why response-decay fits use a predeclared RMS,
upper, or demodulated envelope and exclude neighborhoods of cancellation
zeros.

Finally, the two relevant limits are noncommuting limits:
\label{rem:noncommuting-limits}
\begin{align}
  &\varepsilon\to0\ \text{before the ultimate-time limit}
    &&\Longrightarrow&&
    \text{a Morse--Bott ring and }t^{-1/2},                         \notag\\
  &t\to\infty\ \text{at fixed }\varepsilon\ne0
    &&\Longrightarrow&&
    \text{isolated exact Morse points and }t^{-1}.
  \label{eq:noncommuting-green-limits}
\end{align}
Indeed, \(t_c=|\varepsilon V_4|^{-1}\to\infty\) as
\(\varepsilon\to0\), whereas the coefficient of the separated Morse sum
scales as \(|\varepsilon V_4|^{-1/2}\) and becomes nonuniform when the
points merge into a ring.  The uniform \(J_0\) theorem resolves this
coalescence for \(\varepsilon t=O(1)\); it cannot be extrapolated to
arbitrarily large time at fixed nonzero \(\varepsilon\), because then
\(\varepsilon^2tW_2^{\mathrm{eff}}\) eventually becomes order one.

The results of this section require a finite-radius positive-frequency
ring, nonzero radial curvature, a simple smooth branch with a uniform local
eigengap, a stationary set separated from the window boundary, and exact
nondegenerate Hessians in the fixed-anisotropy theorem.  If \(a=0\), if
\(V_4=0\) without a resolved higher harmonic, if the eigengap closes, or if
an exact Hessian is singular, a different canonical scaling is required.
Material damping introduces an additional exponential envelope and lies
outside the present lossless power-law theorem.  These boundaries prevent
either Eq.~\eqref{eq:uniform-response} or
Eq.~\eqref{eq:morse-stationary-sum} from being applied beyond its proved
regime.
